\section{Integral arithmetic cycles with prescribed residue values}
\label{sa-section}

We now construct the arithmetic cycles used in the class
count.  This proves the source assertion \(\mathrm{SClass}\).
The proof retains the integral ideal lattices, the exact
determinants, and the inverse Galois action in the defining
polynomial.  The matrix size is allowed to depend on the
input; the subsequent stable homology argument is where
that dependence must be removed.

\subsection{A cycle in a lattice stabilizer}

For an element \(g\) of a group, put
\(f_g(a)=g^{-1}ag\).  On bar chains define
\[
 P_g[a_1|\cdots|a_m]
  =\sum_{i=0}^m(-1)^i
    [a_1|\cdots|a_i|g|f_g(a_{i+1})|\cdots|f_g(a_m)].
\]
Here and below the coefficient ring is arbitrary and
commutative.

\begin{lemma}\label{sa-prisms}
The following chain identities hold:
\[
       \partial P_g+P_g\partial=(f_g)_*-1,
\]
and
\[
       \partial Q_{g,h}-Q_{g,h}\partial
                 =P_g+P_h(f_g)_*-P_{gh},
\]
where the degree-two homotopy is
\begin{align*}
 Q_{g,h}[a_1|\cdots|a_m]
   =\sum_{0\le i\le j\le m}(-1)^{i+j}
   [&a_1|\cdots|a_i|g|
       f_g(a_{i+1})|\cdots|f_g(a_j)|h|\\
    &f_{gh}(a_{j+1})|\cdots|f_{gh}(a_m)].
\end{align*}
Consequently, on cycles,
\(P_{gh}-P_g-P_h(f_g)_*\) takes values in boundaries.
\end{lemma}
\begin{proof}
Expand the bar boundaries.  In the first identity,
the faces away from the inserted \(g\) give
\(-P_g\partial\).  The adjacent faces cancel using
\(a g=g f_g(a)\).  The two end faces give
\((f_g)_*\) and minus the identity.

For the second identity, the faces away from both
inserted entries give \(Q_{g,h}\partial\).
Faces next to \(g\) or \(h\) cancel between successive
values of \(i\) or \(j\), using
\(a g=g f_g(a)\) and
\(f_g(a)h=h f_{gh}(a)\).
Deleting an initial \(g\) gives \(P_h(f_g)_*\),
and deleting a final \(h\) gives \(P_g\).
When \(i=j\), merging the adjacent \(g,h\) gives
minus \(P_{gh}\): its sign is
\((-1)^{2i}(-1)^{i+1}=-(-1)^i\).
This accounts for all remaining faces, and proves
the displayed identity with its stated sign.
\end{proof}

\begin{lemma}\label{sa-stabilizer-cycle}
Let \(H\subset G\) be a subgroup, let
\(\rho:D\to G\) and \(d_t:C\to D\) be homomorphisms,
and let \(z\in C_m(C;A)\) be a cycle.
Suppose
\[
                 \partial b=\sum_t n_t(d_t)_*z
\]
for a chain \(b\in C_{m+1}(D;A)\).
Let \(g_t\in G\), and fix an index \(t_0\), writing
\(g_0=g_{t_0}\).
Assume that \(g_t\) commutes with
\(\rho(d_t(C))\), that both \(\rho(D)\) and
\(g_0^{-1}\rho(D)g_0\) lie in \(H\), and that
\(g_0^{-1}g_t\in H\) for every \(t\).
Then there is a cycle \(y\in C_{m+1}(H;A)\)
whose image in \(G\) is homologous to
\[
                    \sum_t n_tP_{g_t}(\rho d_t)_*z.
\]
\end{lemma}
\begin{proof}
Regard \(\rho_0=\rho\) and
\(\rho_1=g_0^{-1}\rho g_0\) as homomorphisms into
\(H\), and put \(k_t=g_0^{-1}g_t\in H\).
Commutation gives
\(f_{k_t}\rho_1d_t=\rho_0d_t\).
Set
\[
 y=(\rho_1)_*b-(\rho_0)_*b
                +\sum_t n_tP_{k_t}(\rho_1d_t)_*z.
\]
By the first identity of Lemma~\ref{sa-prisms}, its
boundary is
\[
 ((\rho_1)_*-(\rho_0)_*)\partial b
       +\sum_t n_t((\rho_0d_t)_*-(\rho_1d_t)_*)z=0.
\]
In \(G\), the second identity of that lemma gives
\[
 \sum_t n_tP_{g_t}(\rho d_t)_*z
 \equiv P_{g_0}\rho_*\partial b
                 +\sum_t n_tP_{k_t}(\rho_1d_t)_*z
                    \pmod{\text{boundaries}}.
\]
The first identity, now applied to \(\rho_*b\),
identifies its first term modulo boundaries with
\((\rho_1)_*b-(\rho_0)_*b\).  The result is \(y\).
\end{proof}

\subsection{Integral modules and good reductions}

\begin{lemma}\label{sa-projective}
Every finitely generated torsion-free module over a
Dedekind domain is a direct summand of a finite free
module.
\end{lemma}
\begin{proof}
For a nonzero fractional ideal \(I\), choose finite
families \(u_i\in I\), \(v_i\in I^{-1}\) with
\(\sum_i u_iv_i=1\), using \(II^{-1}=R\).
The maps
\[
 R^r\longrightarrow I,\quad (a_i)\longmapsto\sum_i a_iu_i,
 \qquad
 I\longrightarrow R^r,\quad x\longmapsto(v_ix)_i
\]
have composite the identity on \(I\).
Thus every such ideal is a finite projective module.

For a finitely generated torsion-free module \(M\)
of positive rank, choose a nonzero linear functional
on \(M\otimes_R\operatorname{Frac}(R)\).
Its image on \(M\) is a nonzero fractional ideal
\(I\).  The surjection \(M\to I\) splits by the
preceding construction.  Its kernel is finitely
generated, torsion-free, and has one less rank;
finite generation follows because \(R\) is Noetherian.
Induction expresses \(M\) as a finite direct sum
of fractional ideals.  Taking the direct sum of
their finite free splittings proves the assertion.
Rank zero gives the zero module, since \(M\) is
torsion-free.
\end{proof}

Let \(F/L\) be a finite Galois extension, and put
\(R=\mathcal O_L\), \(S=\mathcal O_F\).
The module \(S\) is finite over \(R\): an integral
basis over \(\mathbb Z\) already generates it over
\(R\).  It is torsion-free, and hence projective
by Lemma~\ref{sa-projective}.

\begin{lemma}\label{sa-split-reduction}
There exists \(0\ne b\in R\) such that, for every
field \(k\) and homomorphism \(r:S\to k\) whose
restriction to \(R\) is surjective and for which
\(r(b)\ne0\), the map
\[
 S\otimes_R k\longrightarrow
              \prod_{\tau\in\operatorname{Gal}(F/L)} k,
 \qquad x\otimes a\longmapsto(a\,r(\tau x))_\tau
\]
is an isomorphism.
We may enlarge \(b\) so that any specified finite
list of elements of \(F^\times\) are units in
\(S[1/b]\), and any specified finite list of
\(L\)-linear automorphisms of \(F^J\) preserve
\(S[1/b]^J\), for a finite index set \(J\).
\end{lemma}
\begin{proof}
Choose an integral \(\vartheta\in S\) with
\(F=L(\vartheta)\).  Such an element can be obtained
without a separability assumption beyond the one
already present: take an integral basis
\(a_1,\ldots,a_n\) of \(S\) over \(\mathbb Z\)
and consider \(\sum_i t^{i-1}a_i\).
For each nonidentity \(\tau\in\operatorname{Gal}(F/L)\),
the equality between this element and its \(\tau\)-image
is a nonzero polynomial equation in the integer \(t\).
Avoiding the finitely many roots for all \(\tau\)
makes the orbit have \([F:L]\) elements, as required.

Let \(f(T)\in R[T]\) be the monic minimal polynomial
of \(\vartheta\).  Its coefficients are integral
because they are symmetric functions of the integral
conjugates, and they lie in \(L\).
Its discriminant \(\delta\in R\) is nonzero.
As \(R[\vartheta]\) spans \(F\) over \(L\),
clearing the denominators of a finite set of generators
of \(S\) gives \(0\ne a\in R\) such that
\(S[1/a]=R[1/a][\vartheta]\).
Take initially \(b=a\delta\).

In a good residue field, the elements
\(r(\tau\vartheta)\) are all distinct: the product
of their squared differences is the nonzero residue
of \(\delta\).
They give all the roots of \(f\) in \(k\).
Chinese remainders applied to these distinct linear
factors give
\(k[T]/(\overline f)\simeq\prod_\tau k\), with
the stated evaluation maps.  This is the required
tensor-product isomorphism.

For the extra conditions, clear denominators of
each specified nonzero element and its inverse.
For an automorphism, clear the denominators of
the images of a finite set of \(R\)-generators of
\(S^J\), both for it and for its inverse.
Multiplying \(b\) by these finitely many nonzero
elements of \(R\) proves all assertions together.
\end{proof}

If a group \(H\) acts on the projective module
\(\mathcal M=S^J\), choose a complement \(P\) and
an isomorphism \(\mathcal M\oplus P\simeq R^N\).
Let \(H\) act trivially on \(P\).
This embeds the representation in \(\operatorname{GL}_N(R)\).
After reduction, \(\mathcal M\otimes_R k\) is the
space in Lemma~\ref{sa-split-reduction}, repeated
over \(J\).
By Corollary~\ref{sf-primitive-family}, any positive-degree
cycle has the same value under an additive family
before and after the trivial complementary summand
is added.

\subsection{The exact arithmetic construction}

Fix a prime \(\ell\), integers \(\nu\ge c\ge1\),
and assume \(c\ge2\) if \(\ell=2\).
Put
\[
 q=\ell^\nu,\qquad \kappa=1+\ell^c,
 \qquad d=\ell^{\nu-c},\qquad A=\mathbb Z/q\mathbb Z.
\]
Suppose \(F/L\) is cyclic of degree \(d\), generated
by \(\gamma\), and that \(F\) contains a primitive
\(q\)-th root \(\zeta\) with
\(\gamma\zeta=\zeta^\kappa\).
The order of \(\kappa\) in \(A^\times\) is exactly
\(d\).  This follows from the valuation identity
\(v_\ell(\kappa^m-1)=c+v_\ell(m)\) proved in
Lemma~\ref{sc-vandermonde}.

Let \(E\ge1\), assume \(M:=2E+1<\ell^c\), and put
\(a=\kappa^E\).
For \(\alpha\in F^\times\), suppose an integral
divisor \(\mathfrak d\) satisfies
\[
       v_P(\alpha)=\mathfrak d(\gamma P)
                                     -a\mathfrak d(P)
                   \qquad\text{for every prime }P.
\]
Define the weighted norm and the integer factor
\[
 \mathcal N(\alpha)=\prod_{j=0}^{d-1}(\gamma^j\alpha)^{a^j},
 \qquad
 c_E=\prod_{\substack{0\le j<M\\j\ne E}}(a-\kappa^j).
\]

\begin{theorem}\label{sa-class}
There are an integer \(N\), a cycle
\(y\in C_{2E+1}(\operatorname{GL}_N(R);A)\),
and \(0\ne b\in R\) with the following property.
Let \(r:S\to k\) be a homomorphism to a finite field,
surjective on \(R\), with \(r(b)\ne0\).
Let \(v:k^\times\to(A,+)\) be a homomorphism, and
let \(\theta\) be an additive family which vanishes
on boundaries.  For \(x\in A\), write
\(\rho_x(j)=\overline\zeta^{\,xj}I_V\) for the scalar
representation of \(C_q\) on \(V\), and suppose
\[
       \theta_V\bigl(P_g(\rho_x)_*z_E\bigr)
                             =x^E v(\det g).
\]
Then \(\mathcal N(\alpha)\) is a unit in \(S[1/b]\),
and
\[
       \theta_{k^N}(\operatorname{red}(y))
                       =c_E\,v\bigl(r(\mathcal N(\alpha))\bigr)
                        \qquad\text{in }A.
\]
Here reduction of a localized unit means its well-defined
nonzero residue.  The family constructed in
Theorem~\ref{sf-functional} has the required properties
whenever \(\overline\zeta\) remains primitive.
\end{theorem}
\begin{proof}
We give the construction in five steps.

\emph{The cyclic relation and ideal exponents.}
Let \(J=A^2\), and let
\(d_t:C_q\to C_q^2\) be \(j\mapsto(j,tj)\),
for \(t=0,\ldots,M\).
Corollary~\ref{sh-diagonal-relation} gives coefficients
\(n_t\in A\), an index \(t_0\) with \(n_{t_0}\)
a unit, and a chain \(b_1\) such that
\[
                 \partial b_1=\sum_{t=0}^M n_t(d_t)_*z_E.
\]
Write
\[
 P(T)=\prod_{j=0}^{M-1}(T-\kappa^j)=\sum_{i=0}^M p_iT^i,
 \qquad Q(T)=\frac{P(T)}{T-a}\in\mathbb Z[T].
\]
Theorem~\ref{se-interpolation} gives integer exponents
\(\eta_j\), \(j\in J\), with zero pushforward under
\(\chi_t(j_1,j_2)=j_1+tj_2\) for \(t\ne t_0\),
and pushforward \(\sum_i p_i[\kappa^i]\) for \(t=t_0\).
Let \(I=\prod_P P^{\mathfrak d(P)}\), and put
\[
        V=F^J,\qquad \mathcal M=S^J,
        \qquad \mathcal M'=\bigoplus_{j\in J}I^{\eta_j}.
\]
On \(V\), let \(C_q^2\) act diagonally by
\[
     \rho(u_1,u_2)(x_j)_j
            =(\zeta^{j_1u_1+j_2u_2}x_j)_j.
\]
Both ideal lattices are preserved by this action.

\emph{Isomorphisms commuting with the characters.}
For \(t\ne t_0\), divide the coordinates into the
blocks \(\chi_t(j)=x\), \(x\in A\).
The product of the target ideals in each block is
\(I^{\sum_{\chi_t(j)=x}\eta_j}=S\).
Theorem~\ref{si-prescribed} gives an \(F\)-linear
isomorphism in each block, of determinant \(1\),
from its standard lattice to its target lattice.
Their direct sum \(g_t\) maps \(\mathcal M\) onto
\(\mathcal M'\) and commutes with \(\rho d_t\).
Empty blocks have determinant \(1\), as required.

For the distinguished slope, write \(\chi=\chi_{t_0}\).
For \(x=\kappa^i\in\langle\kappa\rangle\subset A^\times\),
put \(T_x=\gamma^i\) and \(B(x)=1\).
For all other \(x\), put \(T_x=1\) and \(B(x)=x\).
The equality of the orders of \(\kappa\) and \(\gamma\)
makes \(T_x\) well defined.
The coordinatewise \(L\)-linear isomorphism
\[
             U(x_j)_j=(T_{\chi(j)}^{-1}x_j)_j
\]
preserves \(\mathcal M\) and changes \(\rho d_{t_0}\)
to the diagonal character with weight \(B(\chi(j))\).
The target ideal in its \(j\)-th coordinate has
divisor
\[
             \eta_j\,\mathfrak d(T_{\chi(j)}P).
\]

For a polynomial \(Q(T)=\sum_i q_iT^i\), define
\[
                       \beta=\prod_i(\gamma^{-i}\alpha)^{q_i}.
\]
The divisor hypothesis on \(\alpha\) gives
\begin{align*}
 v_P(\beta)
 &=\sum_i q_i\bigl(\mathfrak d(\gamma^{i+1}P)
                                      -a\mathfrak d(\gamma^iP)\bigr)\\
 &=\sum_i p_i\mathfrak d(\gamma^iP).
\end{align*}
Thus, in the transformed block of weight \(1\),
the product of the target ideals is \((\beta)\).
Indeed this follows by pushing forward \(\eta\)
and using \(T_{\kappa^i}=\gamma^i\), with repeated
residue classes grouped together.
In every other block the product ideal is \(S\):
there \(T_x=1\), and the relevant sum of exponents
is zero.
The weight-one block is nonempty, since
\(j=(1,0)\) has \(\chi(j)=1\).
Theorem~\ref{si-prescribed} therefore gives a block
map \(g'\) with determinant \(\beta\) in this block
and determinant \(1\) in all the others, carrying
\(\mathcal M\) onto \(U\mathcal M'\).
Put \(g_{t_0}=U^{-1}g'U\).
It maps \(\mathcal M\) onto \(\mathcal M'\) and
commutes with \(\rho d_{t_0}\).
It is only required to be \(L\)-linear; the
coordinatewise Galois changes need not be \(F\)-linear.

\emph{The integral cycle.}
Let \(H=\operatorname{Aut}_R(\mathcal M)\).
All \(g_t\) have the same image lattice, so
\(g_{t_0}^{-1}g_t\in H\).
Since \(\rho\) preserves both lattices, both
\(\rho\) and \(g_{t_0}^{-1}\rho g_{t_0}\) take
values in \(H\).
Choose \(b\) as in Lemma~\ref{sa-split-reduction},
large enough that every \(g_t\) preserves
\(\mathcal M[1/b]\) and that \(\alpha\) and
\(\beta\) are units in \(S[1/b]\).
We may include \(q\) as a factor of \(b\).
Apply Lemma~\ref{sa-stabilizer-cycle} inside
\(\operatorname{Aut}_{R[1/b]}(\mathcal M[1/b])\).
It gives a cycle \(y_0\) in \(H\) homologous,
in this localized group, to
\[
                       \sum_t n_tP_{g_t}(\rho d_t)_*z_E.
\]
Divide \(y_0\) by the unit \(n_{t_0}\in A\),
and extend the representation of \(H\) by a trivial
complement as described after Lemma~\ref{sa-split-reduction}.
Its image is the required integral cycle
\(y\in C_{2E+1}(\operatorname{GL}_N(R);A)\).

\emph{Evaluation after reduction.}
Fix a good residue map \(r\).
All the preceding homology relations can be reduced,
because their chains lie in the localized group.
The trivial complementary summand has zero value.
Under the splitting in Lemma~\ref{sa-split-reduction},
the diagonal character on the component indexed by
\(\tau=\gamma^j\) has its weights multiplied by
\(\kappa^j\) modulo \(q\).
The coordinatewise Galois change \(U\) is an
invertible linear change after reduction; explicitly,
it permutes the embedding coordinate by
\(\tau\mapsto\tau T_{\chi(j)}^{-1}\) in coordinate \(j\).
Therefore invariance under isomorphisms and additivity
allow evaluation in the blocks used above.

Each block prism is a cycle: its scalar cyclic
character commutes with its block map, so the first
identity of Lemma~\ref{sa-prisms} and
\(\partial z_E=0\) apply.
For \(t\ne t_0\), every block determinant is \(1\),
and its value is zero.
For \(t=t_0\), only the weight-one block contributes.
Its determinant on the \(\gamma^j\)-component is
\(r(\gamma^j\beta)\).
The normalization assumed for \(\theta\) gives
\[
       \theta\bigl(P_{g_{t_0}}(\rho d_{t_0})_*z_E\bigr)
               =\sum_{j=0}^{d-1}a^j v(r(\gamma^j\beta)).
\]
The chosen division by \(n_{t_0}\) leaves exactly
this value for \(y\).

\emph{The polynomial and the weighted norm.}
On \(S[1/b]^\times\), put \(f(x)=v(r(x))\), and
\[
                         F(x)=\sum_{j=0}^{d-1}a^j f(\gamma^jx).
\]
The base element \(b\) is fixed by the Galois group,
so all these arguments are localized units.
Since \(a^d\equiv1\pmod q\), reindexing the cyclic
sum gives
\(F(\gamma^{-1}x)=aF(x)\) in \(A\).
Also \(F(xy)=F(x)+F(y)\), including negative powers.
Consequently
\[
     F(\beta)=Q(a)F(\alpha)
            =c_E\,f(\mathcal N(\alpha)).
\]
This is precisely the required value, and
\(\mathcal N(\alpha)\) is a localized unit because
every factor defining it is one.
\end{proof}

All the hypotheses and the evaluation constant in
Theorem~\ref{sa-class} agree with \(\mathrm{SClass}\),
after writing \(E=e+1\).
The proof makes no principal-ideal assumption,
does not confuse an \(L\)-linear map with an
\(F\)-linear map, and retains the complementary
module when embedding in a free integral module.
The construction supplies the arithmetic classes.
Proposition~\ref{sb-odd-rank} and Theorem~\ref{sb-upper}
then prove the stable rational homology estimate
\(\mathrm{BOdd}\) and the finite-image bound
\(\mathrm{SUpper}\); Theorem~\ref{sw-few} combines them
with these classes to complete the finite-set comparison.


